\chapter{死灵夜宴}

“莉莉姆，醒醒。莉莉姆！”莉莉姆感到有人在摇动她的身体。

“嗯？床单怪兽吗……别打扰本小姐睡觉……”莉莉姆不想离开美妙的梦境。

“莉莉姆！”莉莉姆感到温暖的触感；有人捧起了自己的脸，挤压着两侧的脸颊。在肌肉群的带动下，她的眼睛张开了，看到了面前玲的双眼。“怎么了，玲？”莉莉姆张开自己被嘟起的嘴。

“它来了。”玲严肃的表情下隐藏着紧张。莉莉姆的面容凝固了。

玲拉起莉莉姆的手跑出营地。“别出来哦！让姐姐们去看看。”莉莉姆回头提醒睡眼惺忪的雪莉。

云层已经散去。繁星正在哭泣。它们的眼泪垂下天穹，汇聚成流，淌在残破的大地上。乌黑的原浆与这来自天际的洪水相接，或被同化为银镜般的颜色，或被消灭为无声的呜鸣。整座城市的废墟为之震颤，被这冰冷的温暖吞噬。

“它正在吞噬着原浆……”

“它正在吞噬着世界。”莉莉姆指正。

玲看了看手中的移液枪，又看了看莉莉姆手中的长枪。她们仅有的武器如何与尺度相差数个数量级的敌人抗争？

“回去吧，玲。”莉莉姆带着僵死的表情，“这个时代就要结束了；世界决定开启下一个时代。”

“回哪里去？回去做什么？”玲不肯放开莉莉姆的手。

“回营地去。至少最后我们还可以在一起。”莉莉姆挤出不能算是笑容的笑容，“本小姐还没睡饱呢！”

“你在说什么？”玲完全无法理解莉莉姆此刻突然改变的性情，“不如你先告诉我：这究竟是什么？”

“我早该料到的……”莉莉姆撇过头，不想再看那刺眼的反光，“在我们击败其余所有的死灵之后，这就会是最后的敌人，”莉莉姆苦笑着，“它是拟似理智直观。”

流动的镜面以一种无法用语言描述的方式扩散着，遮蔽着、吞噬着它所及的一切。它不像洪水那样受着重力支配，而是不经中介地附着到任何地方，消解了所有地形的差异。这是一种直接的、反感性的、反推论的、反概念的流动方式。它不是通过感官、也不是通过逻辑推理来传达它的到来，而是把对于它的体验直接嵌入潜空间所有居民的脑海。

“这位小姐说的对。”身后传来熟悉的声音。茶会一行人带着他们的八仙桌出现。臭蛤蟆端起茶杯品了品最后的原浆，“我们能够这样看到它，那么它就不是真正的理智直观，而是再生产的拟造物。”臭蛤蟆放下茶杯，“我也不认为达到真正的理智直观是可思议的。”

“注意你的语言。”帽匠还在调整他的帽子。

“不过似乎有人和我的想法不同。”臭蛤蟆叹了口气。

身旁闪过一个奔跑的身影。这个身影踉跄地向着镜子冲去。断壁残垣绊倒了他，但却不能阻止他；他开始手脚并用地在地上飞快地爬行，好似一只大型金毛犬。玲看清了他脸上的镶钻面具，却看不清他面具下的表情；玲认出这是画匠。他在呐喊，他在狂吠：“这就是我！这就是我！这就是我！这就是我！这就是我！这就是我！这就是我！……”

“等一下！那不是！……”玲想阻止他，却是徒劳。

画匠看到了映照在镜子上的他的戴着镶钻面具的面容；玲看清了他的狂妄地笑着的嘴角。他一头扎进这面镜子中，就此消失了。镜子没有产生丝毫裂纹，甚至没有产生任何声响，仿佛根本不在乎他的决绝。

“他不过是绝对者的私生子。”臭蛤蟆评论道。

“以有限逐无限则殆矣。”帽匠评价道。

桌旁的那匹马抱着一个婴儿恸哭。

臭蛤蟆举起蹼向玲和莉莉姆招了招手，想要邀请她们加入这场最后的茶会。玲拒绝了。臭蛤蟆尴尬地笑了笑。他提起茶壶想要为自己续上咖啡，却发现茶壶已经空了。

“一定有办法的！”玲再度攥紧莉莉姆尝试挣脱的手。

“这是世界的意志，”莉莉姆转身离去，“反抗不会成功、也没有意义。”

“等等！”玲的身体被莉莉姆带着向前倾倒。她失去了平衡，跪倒在地。她仍然握着莉莉姆的手腕，“那样的世界绝不是正确的。”

“死脑筋，”莉莉姆不知是生气还是发笑；玲仅仅能看到视线上方她的背影。“到了现在还在纠结正确与否。世界是否正确是由你来定义的吗？”莉莉姆挣开玲的手。

玲同样对这般任性的莉莉姆感到不满。她重新站起来，拼命地追上莉莉姆，想要钳住她的肩膀，却不想反手被莉莉姆摔倒在地，制于身下。玲喘着粗气；她想要平复自己的呼吸，却无法做到，只能凌乱的喊道：“那你呢？现在倒谈起意义来了。难道你之前的战斗仅仅是为了成功，或者说向世界谄媚？”

玲看到莉莉姆的紫色的眼眸湿润了；她嚅动嘴唇，发出几乎不能听闻的声音，“不是‘我的战斗’，而是‘我们的战斗’；你还是和很久以前一样……”

玲尝试说服莉莉姆，“总有办法的！你不是说了潜空间是所有可能性的应许之地吗……”玲停止了话语，因为莉莉姆的眼泪滴进了她的眼眶。玲却没有丝毫想要眨眼以排除异物的冲动，可能是因为她也早已开始哭泣。两个人就这样上下对视，互望泪眼；言语在此刻消失。
